\subsection{The actual central word realizing the quadratic coefficient}
\label{sec:nonstandard-central-word}

\paragraph{P-F52 (an intertwiner inside the displayed nonstandard algebra).}
Use the cubic algebra $B_0$ over $F_0=\mathbb C(q)$, its splitting
field $K=F_0(x)$, and the elements $c,c',P,D,F$ of P-F36 and P-F39.
Here $x^2=D$; the letter $F$ without a subscript continues to denote
the positive Laurent normalization in (P-F39a).
Write $A=Q_1$, $B=Q_2$, and
\[
 f=(q^4+1)(q^4-q^2+1)(q^2+1)/q^5
   =q^5+q+q^{-1}+q^{-5}.
\]
Define the following explicit Laurent polynomial and central word:
\[
 \begin{split}
 S={}&q^{10}+q^{-10}+3(q^8+q^{-8})+(q^6+q^{-6})
       +5(q^4+q^{-4})-2(q^2+q^{-2})+16,\\
 Z={}&(Q_1-Q_2)^2+(S-f^2)1.
 \end{split}                                                    \tag{P-F52a}
\]
In the actual source blocks $U_1,U_2\subset B_0\otimes K$,
left and right multiplication by $Z$ have scalar values $-c,-c'$,
respectively. In particular, this operator is specified by a short
word in the original generators, rather than by assigning arbitrary
eigenvalues to abstract central idempotents.

There is an explicit two-dimensional $F_0$-subspace of $B_0$ on
which multiplication by $Z$ has precisely the matrix $M_0$ of
(P-F39k). Its multiple $FZ$ has matrix $M=FM_0$.
The same construction gives four copies of this matrix on the
eight-dimensional descended two-sided ideal associated with
$U_1\oplus U_2$. These are statements about actual multiplication
in $B_0$, with the original source basis retained.

For every real $q>0$, the same word $Z$ acts as a strictly positive,
self-adjoint operator in the source's actual tensor representation
on $X_q^{\otimes3}$, with the explicit positive tensor form given
below. So does $FZ$. This last assertion concerns the tensor
specialization, whose image can have smaller dimension at exceptional
parameters; the generic $F_0$-linear intertwiners are not asserted to
remain injective at every specialization.

\paragraph{Only the original relation and the original cell element are needed.}
The source gives $A^2=fA$, $B^2=fB$, and the independent
alternating-word basis of lengths at most ten: GCT VII,
Section~7.1.1, pp.~18--19. Put
\[
       C=(A-B)^2=f(A+B)-AB-BA,
       \qquad G=f^2-C,
       \qquad w_j=A(BA)^j.
                                                               \tag{P-F52b}
\]
Direct multiplication gives
$AC=CA=f^2A-ABA$ and $BC=CB=f^2B-BAB$.
Thus $C$ and $G$ are central, and
\[
                    Gw_j=w_{j+1}\qquad(j\ge0).                 \tag{P-F52c}
\]
The actual relation in Figure~2, printed p.~31, has the form
\[
            w_5-Sw_4+a_3w_3+a_2w_2+a_1w_1+a_0w_0=0,          \tag{P-F52d}
\]
where the coefficient of $w_5$ is one and the coefficient of
$w_4$ is exactly $-S$ from (P-F52a). The actual element $u=u_1^1$ in
Figure~7, printed p.~36, has the form
\[
                  u=w_4+cw_3+b_2w_2+b_1w_1+b_0w_0.            \tag{P-F52e}
\]
Its conjugate $u'=u_2^1$ replaces $x$ by $-x$ and therefore has
coefficient $c'$ at $w_3$. This is the same scalar coefficient
used in GCT VIII and P-F39, even though VIII also substitutes
$P$-words for $Q$-words in its corresponding element.

The following full coefficient identity verifies the eigenvector
equation without relying on the simple-block assertion. Set
\[
 \begin{gathered}
 R=-3q^{28}-4q^{24}+2q^{22}-10q^{20}+2q^{18}-4q^{16}-3q^{12},\\
 g_1=\frac{-R-x}{2q^{20}},\qquad
 g_2=\frac{-R+x}{2q^{20}},\\
 g_3=q^4+2+q^{-4},\quad g_4=q^6+2+q^{-6},\quad
 g_5=q^{10}+2+q^{-10},\\
 E_1=g_3+g_4+g_5,\quad E_2=g_3g_4+g_3g_5+g_4g_5,\quad
 E_3=g_3g_4g_5,\\
 L=-R/q^{20},\qquad
 N_0=\frac{(q^4-q^2+1)^2
             (q^{12}+q^{10}+2q^8+2q^4+q^2+1)^2}{q^{16}}.
 \end{gathered}                                                 \tag{P-F52f1}
\]
Expansion gives $R^2-D=4q^{40}N_0$, $g_1+g_2=L$,
$g_1g_2=N_0$, and $L+E_1=S$. All coefficients of the
original relation and the original cell element can be written
as follows; the expressions are direct factorizations of
Figures~2 and 7, including their lower coefficients:
\[
\begin{array}{c|c|c}
 j&[t^j]r(t)&[t^j]p_1(t)\\ \hline
 0&-N_0E_3&g_2E_3\\
 1&N_0E_2+LE_3&-E_3-g_2E_2\\
 2&-N_0E_1-LE_2-E_3&E_2+g_2E_1\\
 3&N_0+LE_1+E_2&-E_1-g_2=c\\
 4&-L-E_1=-S&1\\
 5&1&0
\end{array}                                                     \tag{P-F52f2}
\]
Here $r(G)A=0$ is exactly (P-F52d), and $p_1(G)A=u$
is exactly (P-F52e). For example their constant coefficients
are the factored first rows of the two source figures; the
other rows follow by the same multiplication and collection
of powers of $q$ and $x$. The accompanying exact verifier
transcribes and compares all six relation coefficients and
all five cell coefficients, without interpolation.
Multiplying the polynomials in the table now gives
\[
 \begin{split}
 r(t)&=(t^2-Lt+N_0)(t^3-E_1t^2+E_2t-E_3),\\
 p_1(t)&=(t-g_2)(t^3-E_1t^2+E_2t-E_3),\\
                 (t-g_1)p_1(t)&=r(t).
 \end{split}                                                    \tag{P-F52f3}
\]
Substitution of the central element $G$ proves
$(G-g_1)u=r(G)A=0$ directly. Since $g_1=S+c$,
\[
                 Gu=(S+c)u,\qquad Zu=-cu.                     \tag{P-F52f}
\]
The leading coefficient gives an independent quick check:
in $Gu$ only the reduction of $w_5$ and the term $cw_4$
contribute to $w_4$, with coefficient $S+c$.
Conjugating gives $Zu'=-c'u'$. Both elements are nonzero because
their coefficient at the independent word $w_4$ is one.
Moreover $c-c'=-x/q^{20}\ne0$, so they are distinct; their
belonging to different direct summands also makes them
linearly independent over $K$.

To label the corresponding full blocks, use the split Wedderburn
decomposition in GCT VII, Section~7.1.2, pp.~19--20: $u,u'$
belong to the four-dimensional simple matrix-algebra summands
$U_1,U_2$. The center of a matrix algebra consists of scalars,
as follows by commuting with its diagonal and off-diagonal matrix
units. Thus (P-F52f) gives the same scalar on the full respective
summands, on either side. The eigenvector identity itself has
already been verified independently in (P-F52f1)--(P-F52f3).

This argument does not use the printed two-by-two generator
matrices on p.~20. As printed, both are upper triangular and
would give a common invariant line, inconsistent with the
generic irreducibility claimed there. That display is
quarantined here; no silent correction to it is needed for
(P-F52f), which follows from the complete original relation
and cell-element factorization. The cited split simple-algebra
assertions serve separately to identify the full block labels.

\paragraph{The exact descended intertwiner.}
Let $\tau(x)=-x$ be the nontrivial automorphism of $K/F_0$.
Set
\[
            v_0=u+u',\qquad v_1=x(u-u').                       \tag{P-F52g}
\]
They are fixed by $\tau$ and hence lie in $B_0$. To see this
last assertion explicitly, expand in the $F_0$-word basis:
each fixed coefficient in $K$ belongs to $F_0$.
The coefficient of $w_4$ in $v_0$ is two and in $v_1$ is zero,
while that of $w_3$ in $v_1$ is
$x(c-c')=-D/q^{20}\ne0$. Thus $v_0,v_1$ are independent.

The $F_0$-linear map
\[
       \iota:K\longrightarrow B_0,
           \qquad a\longmapsto au+\tau(a)u'                  \tag{P-F52h}
\]
is injective and sends the basis $(1,x)$ to $(v_0,v_1)$.
Equation (P-F52f) gives
\[
       Z\,\iota(a)=\iota\left(\frac{P+x}{2q^{20}}a\right).
                                                               \tag{P-F52i}
\]
In that basis, actual multiplication by $Z$ is therefore
\[
                 \frac1{2q^{20}}
                       \begin{pmatrix}P&D\\1&P\end{pmatrix}, \tag{P-F52j}
\]
exactly (P-F39k), with no transpose, sign, or scalar discrepancy.
Multiplying $Z$ by the source normalization $F$ gives (P-F39m).
The map $\iota$ intertwines the indicated central action; its
two-dimensional image is not asserted to be a left ideal of $B_0$.

For the actual two-sided ideal, take the four source basis elements
of $U_1$, namely $u_1^1,u_1^{12},u_1^{21},u_1^2$, and their
conjugates in $U_2$. Call them $u_r,u'_r$, respectively.
The eight vectors $u_r+u'_r$, $x(u_r-u'_r)$ are $F_0$-rational
and independent: a relation, projected onto $U_1$ and $U_2$,
has coefficients $a_r+xb_r=0$ and $a_r-xb_r=0$, forcing
$a_r=b_r=0$. Conversely every fixed vector in
$U_1\oplus U_2$ is of this form by writing its $U_1$ coefficients
as $a_r+xb_r$. Their span is therefore the descended two-sided
ideal $W=(U_1\oplus U_2)^\tau$, of dimension eight over $F_0$.
Multiplication by $Z$ acts on it as four copies of (P-F52j).
The $U_i$ here are matrix-algebra summands, of dimension four;
the source's simple left modules $T_i$ have dimension two.
These two roles are not identified.

\paragraph{An explicit positive form for the actual tensor action.}
For $q>0$ set
\[
 [2]=q+q^{-1},\quad [3]=q^2+1+q^{-2},\qquad
 (d_0,d_1,d_2,d_3)
       =\bigl(1,q^{-1}[3],[3][2]^2,q^3[3]^2[2]^2\bigr).
                                                               \tag{P-F52k}
\]
All four entries are positive. Give the source basis
$x_i=f^ix_0$ of $X_q$ the diagonal Gram matrix
$\operatorname{diag}(d_0,d_1,d_2,d_3)$, and use its tensor
product on tensor powers. Here the lowering operator denoted
$f$ in the source's basis notation is distinct from the scalar
Laurent polynomial $f$ in (P-F52a).

The actual rescaled projection matrix $\mathsf P$ in GCT VII,
Figure~1, printed p.~30, is self-adjoint for this tensor form.
The entire check can be organized by $i+j$, since its nonzero
entries preserve that sum. The one-dimensional blocks and the
two-dimensional blocks $i+j=1,5$ have real diagonal entries
and equal opposite off-diagonal entries, with equal Gram weights.
For $i+j=2$, the weights are $(d_2,d_1^2,d_2)$; the only
nontrivial opposite-entry ratio is verified by
\[
 \frac{d_2}{d_1^2}
       =\frac{q^2(q^2+1)^2}{q^4+q^2+1}.
\]
For $i+j=3$, the weights are
$(d_3,d_1d_2,d_1d_2,d_3)$; the displayed outer-to-inner
opposite-entry ratios are $q^{-4}$, and
$d_3/(d_1d_2)=q^4$. Opposite entries between equal-weight
positions are equal. Finally, for $i+j=4$ the weights are
$(d_1d_3,d_2^2,d_1d_3)$, and the required ratio is
\[
 \frac{d_1d_3}{d_2^2}
       =\frac{q^2(q^4+q^2+1)}{(q^2+1)^2}.
\]
Substitution in the entries of Figure~1 gives
$d_id_j\mathsf P_{ij,kl}=d_kd_l\mathsf P_{kl,ij}$
in every block. These identities are Laurent identities,
so also hold at $q=1$, where some off-diagonal entries vanish;
no division by such an entry is being asserted there.
Thus $\mathsf P$ is self-adjoint. The source's complementary
rescaled projection is $\mathsf Q=fI-\mathsf P$ by
equations~(4), (16)--(18), pp.~6--7 and 18, so it too is
self-adjoint. Consequently $Q_1=\mathsf Q\otimes I$ and
$Q_2=I\otimes\mathsf Q$ on $X_q^{\otimes3}$ are
self-adjoint for the displayed positive product form.

Put $t=q^2+q^{-2}=2+v$, so $v\ge0$. A direct expansion gives
\[
 \begin{split}
 S-f^2
 &=3(q^8+q^{-8})-(q^6+q^{-6})
       +3(q^4+q^{-4})-3(q^2+q^{-2})+12\\
 &=3t^4-t^3-9t^2+12
   =3v^4+23v^3+57v^2+48v+16>0.
 \end{split}                                                    \tag{P-F52l}
\]
For every nonzero tensor vector $w$, self-adjointness now gives
\[
 \langle w,Zw\rangle
   =\|(Q_1-Q_2)w\|^2+(S-f^2)\|w\|^2>0.
\]
This proves actual positive-definiteness, without importing a
generic compact-star assertion across an exceptional specialization.
The coefficients of $F$ are positive, so $F(q)>0$ and the same
conclusion holds for $FZ$.\hfill$\square$

\paragraph{What this finite comparison settles.}
The positive quadratic operator of P-F39 is now realized by an
explicit element of the actual nonstandard algebra, with a
source-basis intertwiner and an explicit tensor positive form.
In particular the occurrence of the eigenvalue $-c$ is tied to
the displayed word coefficient of $u_1^1$ through the leading-term
calculation (P-F52d)--(P-F52f); it is not imposed after choosing an
unrelated matrix model. On the descended two-dimensional space,
however, the unscaled determinant still has coefficients $T,N$
from (P-F39c), with their negative Laurent coefficients. Thus the
actual positivity just proved does not imply coefficientwise
positivity. The prescribed scale $F$, the other source coefficients,
their canonical-basis meanings, and any cohomological or theta
intertwiner remain substantive additional data. No uniform positivity
theorem in tensor degree or PH2/SH consequence is claimed here.

Primary source: GCT VII, Section~7.1.1--7.1.2, pp.~18--20,
Figure~1 p.~30, Figure~2 p.~31, Figure~7 p.~36,
\url{https://arxiv.org/pdf/0709.0749}; GCT VIII's matching scalar
coefficient and its normalization are cited and checked in P-F39,
\url{https://arxiv.org/pdf/0709.0751}.

\paragraph{A finite normalization corollary to P-F52.}
Even the combination of positive actual central action and symmetric
Tate--Lefschetz trace and norm does not characterize the source's
prescribed normalization $F$. Indeed put
\[
       H(q)=(q+1+q^{-1})^4
        =19+16(q+q^{-1})+10(q^2+q^{-2})
                    +4(q^3+q^{-3})+(q^4+q^{-4}).               \tag{P-F52m}
\]
For every $q>0$, $HZ$ is an actual strictly positive central
operator on the same tensor space, since $H(q)>0$.
On the descended two-dimensional space its determinant polynomial is
\[
             \det(I+tHZ)=1+H(q)T(q)t+H(q)^2N(q)t^2,             \tag{P-F52n}
\]
where $T,N$ are precisely (P-F39c).
In the symmetric-vector notation $[v]$ of P-F39, direct convolution
gives the complete coefficient certificates
\[
 \begin{split}
 HT=[\alpha^H],\quad
 \alpha^H={}&(390,336,242,160,139,132,128,112,103,88,70,44,23,8,2),\\
 H^2N=[\beta^H],\quad
 \beta^H={}&(169556,164696,152852,139200,127656,118312,109695,\\
 &100480,90558,79776,68121,55832,44025,33832,25905,20000,\\
 &15558,12016,9093,6632,4587,2936,1699,864,377,136,39,8,1).
 \end{split}                                                    \tag{P-F52o}
\]
Every entry is positive, and each list is nonincreasing.
These are exact finite identities: if $h_j$ are the coefficients
in (P-F52m), extended symmetrically and by zero, then
\[
 \alpha_j^H=\sum_i h_i t_{j-i},\qquad
 \beta_j^H=\sum_{i,k}h_ih_kn_{j-i-k},                           \tag{P-F52p}
\]
using the explicitly displayed coefficients $t_j,n_j$ of (P-F39c).
The two coefficient sums at one are $81\cdot44=3564$ and
$81^2\cdot448=2939328$. The telescoping formula (P-F39h)
therefore realizes each as a nonnegative direct sum of the same
Tate-twisted projective-space Lefschetz characters.

Nevertheless $H$ is not divisible by $F$ in
$\mathbb C[q,q^{-1}]$. For a nonzero Laurent polynomial define
its width as its largest exponent minus its smallest exponent.
The extremal coefficients in a product are nonzero products,
so widths add. Here $F$ has width $32$ by (P-F39a), whereas
$H$ has width $8$, excluding $H=FJ$ for any nonzero Laurent
polynomial $J$. In particular these conditions do not even
force the prescribed factor to divide the normalization.

This proves the precise failure of characterization by the stated
finite operator and determinant properties. It does not assert
that $H$ has the source's additional lattice, integral-form, or
canonical-basis compatibility, or that it normalizes all other
coefficients. Those compatibility conditions remain the missing
datum; positivity of this actual action and its two Lefschetz
coefficients alone cannot supply them.\hfill$\square$

\paragraph{The exact spectral-denominator provenance of $F$.}
The additional rule that specifies this source normalization can
itself be identified explicitly. In GCT VII, p.~27,
equations~(32)--(35), the actual two-tensor braid operator
$\widehat R$ has eigenvalues
\[
       q_1=q^{9/2},\qquad q_2=-q^{-3/2},\qquad
       q_3=q^{-11/2},\qquad q_4=-q^{-15/2}.                    \tag{P-F52q}
\]
The positive spectral projector is expressed by Lagrange interpolation:
\[
  P_+=\sum_{i\in\{1,3\}}
       \frac{\prod_{j\ne i}(\widehat R-q_j)}
            {\prod_{j\ne i}(q_i-q_j)}.
\]
The source specifies the product of the two displayed denominators,
not the smallest common denominator after matrix cancellations:
\[
 f_p=\prod_{i\in\{1,3\}}\prod_{j\ne i}(q_i-q_j).
                                                               \tag{P-F52r}
\]
In particular the factors $q_1-q_3$ and $q_3-q_1$ both occur.
Their product is $-q^{-11}(q^{10}-1)^2$. The other four factors
are respectively
\[
 q^{-3/2}(q^6+1),\quad q^{-15/2}(q^{12}+1),\quad
 q^{-11/2}(q^4+1),\quad q^{-15/2}(q^2+1).
\]
Thus their exact product and its ratio to the convenient matrix
normalizer $f=q^{-5}(q^4+1)(q^6+1)$ are
\[
 \begin{split}
 f_p&=-q^{-33}(q^{10}-1)^2(q^6+1)(q^{12}+1)(q^4+1)(q^2+1),\\
 f_p/f&=-q^{-28}(q^{10}-1)^2(q^{12}+1)(q^2+1).
 \end{split}
\]
The source then removes the displayed vanishing factor by defining
$\widehat f_p=-q f_p/((q-1)^2f)$ on printed p.~28.
Consequently the centering used in P-F39 gives exactly
\[
 \begin{split}
 F=q^{11}\widehat f_p
  &=q^{-16}(1+q+\cdots+q^9)^2(1+q^2)(1+q^{12}).
 \end{split}                                                    \tag{P-F52s}
\]
This is a Laurent identity, not an evaluation at a singular
interpolation parameter. The factorization cancels $(q-1)^2$
and extends regularly to $q=1$, where $F(1)=10^2\cdot2\cdot2=400$.
It proves positivity of the coefficients directly. Expanding it
also gives exactly the complete coefficient list in (P-F39a).

Thus the finite provenance of $F$ is fixed once one supplies the
source's particular interpolation-denominator product and its
specified removal and centering factors. The positive central
action alone does not impose that interpolation rule, as the
alternative $H$ proves. This dictionary also does not assert that
the spectral-denominator rule supplies a categorical integral form,
canonical-basis multiplication, or the required geometric
intertwiner for all nonstandard coefficients. Those additional
compatibilities are distinct from the explicit algebraic
provenance just established.\hfill$\square$

\paragraph{An ordinary lambda-operation obstruction for the two coefficient models.}
The separate positive coefficient models cannot be combined into
the first two ordinary exterior-power characters of one graded
vector space, even after replacing $F$ by any common scalar
normalization. More precisely, there is no rational function
$J(q)$ regular and nonzero at $q=1$ and no finite-dimensional
$\mathbb Z$-graded vector space $V$ for which
\[
 \operatorname{ch}_q V=J(q)T(q),\qquad
 \operatorname{ch}_q\bigwedge^2V=J(q)^2N(q).
                                                        \tag{P-F52t}
\]
The same obstruction holds for a virtual graded-vector-space
class with the ordinary special lambda-ring operations.

Indeed evaluation of a graded character at one is its dimension.
Writing $r=44J(1)$ and using $T(1)=44$, $N(1)=448$ from
P-F39, the two equalities in (P-F52t) would give
\[
 r\in\mathbb Z,\quad r\ne0,\qquad
 \frac{r(r-1)}2=448J(1)^2=\frac{28}{121}r^2.
                                                        \tag{P-F52u}
\]
Cancelling $r$ yields $65r=121$, impossible for an integer.
For virtual classes the same calculation applies: their rank
is an integer and the rank of the second lambda operation is
$\binom r2$, as follows by expanding
$(1+t)^{r}$ (including negative integral exponents).
For the prescribed $F$ the numerical discrepancy is particularly
explicit:
\[
 \binom{17600}{2}=154871200\ne71680000=N_1(1).
                                                        \tag{P-F52v}
\]
Thus individual Tate--Lefschetz character realizations of
$T_1$ and $N_1$ do not give this common-object lambda
dictionary.

There is a related obstruction when $q$ is interpreted as
finite-field cardinality and the same Frobenius is required
over all finite extensions. Fix any integer $p>1$. There is
no single finite-dimensional linear operator $A$ over a
characteristic-zero field whose traces satisfy
\[
 \operatorname{Tr}(A^r)=T_1(p^r),\qquad
 \operatorname{Tr}\bigl(\bigwedge^2 A^r\bigr)=N_1(p^r)
                   \quad\hbox{for every }r\ge1.
                                                        \tag{P-F52w}
\]
To prove this, triangularize $A$ over an algebraic closure.
If its diagonal entries are $a_i$, summing
$a_i^r a_j^r$ over $i<j$ gives the elementary identity
\[
 2\operatorname{Tr}\bigl(\bigwedge^2 A^r\bigr)
       =\operatorname{Tr}(A^r)^2-\operatorname{Tr}(A^{2r}).
\]
The proposed traces would therefore make the Laurent
polynomial
\[
       2N_1(q)-T_1(q)^2+T_1(q^2)
                                                        \tag{P-F52x}
\]
vanish at every $q=p^r$. These are infinitely many distinct
nonzero values, so after multiplication by a power of $q$
the resulting polynomial must vanish identically. Evaluating
at one contradicts (P-F52v). No purity or semisimplicity
assumption is used in this contradiction.

These statements exclude the specified ordinary lambda
dictionary and the specified extension-compatible Frobenius
dictionary. A graded trace of an additional endomorphism,
a different relationship between the two coefficient objects,
or a different operation is not identified with an ordinary
graded dimension. In particular the actual central operator
and its determinant remain valid; the obstruction concerns
their proposed common-object cohomological interpretation,
not the algebraic determinant identity.\hfill$\square$

\paragraph{An explicit positive multiplication table for the actual central algebra.}
There is a further finite positivity statement for the prescribed
normalization itself. It concerns the algebra of restricted central
endomorphisms; the intertwining subspace $\iota(K)\subset B_0$ is
not being asserted to be a subalgebra.

Write $R_0=\mathbb Z[q,q^{-1}]$, $t=T_1$, $n=N_1$, and set
\[
 \begin{split}
 e(q)&=q^4+q^{-4}+q^8+q^{-8}
                    +q^{12}+q^{-12}+q^{24}+q^{-24},\\
 a(q)&=\frac{t(q)-e(q)}2,\qquad
 d(q)=\frac{t(q)+e(q)}2,\\
 b(q)&=a(q)d(q)-n(q)
       =\frac{t(q)^2-4n(q)-e(q)^2}{4}.
 \end{split}                                                    \tag{P-F52y}
\]
Then $a,d,b$ belong to $\mathbb Z_{\ge0}[q,q^{-1}]$.
In fact all three have positive symmetric coefficient vectors,
nonincreasing from the middle towards either end. Here is a
complete finite certificate in the notation $[v]$ of P-F39:
\[
 \begin{split}
 a=[\alpha],\qquad \alpha={}&
 (354,354,347,340,336,333,321,309,282,256,232,208,182,\\
 &157,131,105,85,65,53,41,31,21,15,9,5,2,1),\\
 d={}&a+e,\\
 b=[\beta],\qquad \beta={}&
 (173758,173280,171874,169576,166431,162488,157785,\\
 &152364,146308,139703,132653,125265,117614,109779,101839,\\
 &93877,85986,78264,70768,63560,56676,50156,44035,38352,\\
 &33124,28369,24074,20227,16813,13819,11223,9005,7131,5569,\\
 &4281,3231,2390,1731,1225,845,567,367,228,133,72,35,15,5,1).
 \end{split}                                                    \tag{P-F52z}
\]
The vector for $d$ is obtained by adding one to entries
$4,8,12,24$ of $\alpha$, leaving every other entry unchanged.
These assertions are exact integer Laurent identities. To verify
the certificate directly, extend the displayed vectors
symmetrically and by zero, and write $e_j=[q^j]e$ and
$n_j=[q^j]N_1$, the latter given in P-F39. The two coefficient
identities to check are precisely
\[
 [q^j]T_1=2\alpha_j+e_j,\qquad
 \beta_j=\sum_i\alpha_i(\alpha_{j-i}+e_{j-i})-n_j.
                                                               \tag{P-F52aa}
\]
They give every coefficient, including the vanishing outside
the displayed ranges. Positivity and monotonicity then follow
from the full lists, with the stated four increments for $d$.
Thus each of $a,d,b$ also has the explicit Tate--Lefschetz
character realization of the telescoping formula (P-F39h).
The structure coefficient $e$ is nonnegative and symmetric;
no unimodality is asserted for $e$.

Let $W_0=\iota(K)$ be the two-dimensional $F_0$-space in
(P-F52h), and let $\mathsf z$ be the restriction to $W_0$
of actual multiplication by $FZ$. Define
\[
 \mathcal A=F_0[\mathsf z]\subset\operatorname{End}_{F_0}(W_0),
 \qquad \mathsf u=\mathsf z-a\,\mathrm{id}_{W_0}.
                                                               \tag{P-F52ab}
\]
Under the existing intertwiner, $\mathsf z$ is multiplication
in $K$ by $F(P+x)/(2q^{20})=-Fc$. This element still generates
$K/F_0$, since $F\ne0$ and $c$ has degree two by P-F36.
Consequently $\mathcal A\cong K$, and its minimal polynomial
is $X^2-tX+n$. Substituting $\mathsf z=a+\mathsf u$,
using $t=a+d$ and $d-a=e$, gives the positive multiplication
table
\[
       \mathsf u^2=e\mathsf u+b\,\mathrm{id},\qquad
 \mathcal O:=R_0[\mathsf z]=R_0[\mathsf u]
       =R_0\mathrm{id}\oplus R_0\mathsf u.
                                                               \tag{P-F52ac}
\]
The last sum is direct because $1,\mathsf u$ are independent
over $F_0$. Polynomial division by the monic quadratic shows
that
\[
           \mathcal O\cong R_0[U]/(U^2-eU-b).
\]
Indeed every class has a unique candidate representative of
degree at most one, and independence over $F_0$ makes that
representative unique in the actual endomorphism algebra.
More precisely, $\mathbb Q(q)[\mathsf z]$ is already a
quadratic field, since the minimal polynomial is irreducible
even over $F_0=\mathbb C(q)$. The algebra $\mathcal O$ is a
free $R_0$-order in that field, and its extension of scalars
to $F_0$ is the actual algebra $\mathcal A$.
Its displayed basis has nonnegative integer Laurent structure constants.
It is not a claim about the source's full integral form or
its conjectural canonical basis.

The regular action of this order is also identified with the
actual central action on the original descended space. Put
\[
       w=\frac{e}{2}v_0+\frac{F}{2q^{20}}v_1,
       \qquad v_0,v_1\text{ as in (P-F52g)}.
                                                               \tag{P-F52ad}
\]
Then $w=\iota(F(P+x)/(2q^{20})-a)$ by $t=FP/q^{20}$.
The change from $(v_0,v_1)$ to $(v_0,w)$ has determinant
$F/(2q^{20})\ne0$, so it is a basis over $F_0$.
The map $\mathcal A\to W_0$ sending
$r(\mathsf z)$ to $r(\mathsf z)v_0$ sends $(1,\mathsf u)$
to $(v_0,w)$ and is an isomorphism of $\mathcal A$-modules.
Thus actual multiplication by $FZ$ has the matrix
\[
             \mathsf W(q)=
                 \begin{pmatrix}a(q)&b(q)\\1&d(q)\end{pmatrix}
                    \quad\text{in the basis }(v_0,w).
                                                               \tag{P-F52ae}
\]
This proves an explicit similarity to the matrix $FM_0$ of
P-F39, with the prescribed source normalization and without
assigning new abstract eigenvalues. The same change can be
made in each of the four copies on the descended ideal of
(P-F52j). As before these are generic $F_0$-linear statements;
specialized injectivity at every value of $q$ is not asserted.

For every integer $r\ge1$ it follows in particular that
\[
 \operatorname{Tr}_{W_0}((FZ)^r)
       =\operatorname{Tr}(\mathsf W(q)^r)
       \in\mathbb Z_{\ge0}[q,q^{-1}].
                                                               \tag{P-F52af}
\]
There is a literal finite counting formula for each coefficient.
Take a directed graph on vertices $1,2$, with
$[q^j]\mathsf W_{ki}$ distinguishable edges of degree $j$
from $i$ to $k$. For a fixed $r$, matrix multiplication
counts the length-$r$ walks with degree equal to the sum of
edge degrees: this follows inductively from its sum over the
intermediate vertex and the convolution of Laurent
coefficients. The trace counts closed walks with a specified
starting vertex. Negative edge degrees are permitted, and
every fixed-length count is finite. This is not the identity
$\operatorname{Tr}(\mathsf W(q)^r)=T_1(q^r)$, which is
neither needed nor asserted; the extension-compatible
Frobenius obstruction above is therefore unchanged.

Finally, positivity of this table does not supply a rigid
simple-object categorification. There is a concrete obstruction
to that particular interpretation. At $q=1$,
\[
             e(1)=8,\qquad b(1)=5759984,
             \qquad u^2=8u+5759984\,1.
                                                               \tag{P-F52ag}
\]
Suppose a complex-linear semisimple rigid tensor category
with finite-dimensional morphism spaces and simple unit
had Grothendieck ring identified with $\mathcal O/(q-1)$,
with its simple classes $[\mathbf1],[U]$ corresponding to
the displayed basis $1,u$.
The unit multiplicity in $U\otimes U$ would be $5759984$.
Rigidity instead identifies
$\operatorname{Hom}(\mathbf1,U\otimes U)$ with
$\operatorname{Hom}(U^*,U)$ for a suitable dual.
The dual of a simple object is simple, and Schur's lemma
makes this dimension at most one. This is a contradiction.
The assertion concerns this simple-basis interpretation
after $q=1$, not all possible models involving nonsimple
objects, different operations, or different bases.

There is, however, an explicit additive character-level model
of the central operator. Let $\mathscr V$ be the category of
finite-dimensional $\mathbb Z$-graded complex vector spaces.
Choose the finite Tate--Lefschetz coefficient spaces
$A_{\rm coh},B_{\rm coh},D_{\rm coh}$ just specified by
(P-F39h) and (P-F52z), with characters $a,b,d$, respectively.
Define an additive endofunctor of $\mathscr V\times\mathscr V$ by
\[
 \mathscr F(V_1,V_2)=
 \bigl(A_{\rm coh}\otimes V_1\ \oplus\ B_{\rm coh}\otimes V_2,
                    \ V_1\ \oplus\ D_{\rm coh}\otimes V_2\bigr).
                                                               \tag{P-F52ah}
\]
Morphisms are sent by the corresponding tensor and direct-sum
operations. The functor commutes with grading shifts.
Every finite graded vector space is a finite direct sum of
one-dimensional shifts, so
$K_0(\mathscr V\times\mathscr V)=R_0^2$.
Tensor product multiplies its graded characters, while
direct sum adds them. Hence $K_0(\mathscr F)$ has exactly
the matrix $\mathsf W(q)$, not merely the same eigenvalues.
After extension of scalars to $F_0$, the map sending the
two standard basis vectors to $v_0,w$ is an isomorphism
and intertwines $K_0(\mathscr F)$ with the actual restricted
central action of $FZ$ by (P-F52ae).
This supplies a finite cohomological coefficient model and an
exact comparison on its Grothendieck group. It does not
identify these large coefficient spaces with $W_0$ itself,
make $\mathscr F$ a monoidal functor, or realize the full
nonstandard algebra and its canonical basis categorically.

The finite construction therefore supplies actual positive
Laurent multiplication and all central power-trace counts,
while the source-canonical integral form, the other
nonstandard generators and products, coherent monoidal
data, and a uniform construction in tensor degree remain
additional requirements.\hfill$\square$

\paragraph{Compatibility with the actual word form and internal bar.}
The preceding positive lattice has more source compatibility than
an abstract integral order, but the normalization of the generators
matters. In this paragraph the ``word form'' means exactly
\[
       \mathcal B_{\rm word}
           =R_0\langle Q_1,Q_2\rangle\subset B_0,
           \qquad R_0=\mathbb Z[q,q^{-1}],
                                                               \tag{P-F52ai}
\]
with the convenient generators of GCT VII, equations~(16)--(18),
printed p.~18. These are also the generators explicitly used in
GCT VIII, Section~5.3, p.~31. Since $P_i=f-Q_i$ and $f\in R_0$,
the same form is generated by the convenient $P_i$.

This form is free on the source's $21$ alternating words of
length at most ten, including the empty word. Indeed the
complete presentation in GCT VII, Section~7.1.1, pp.~18--19,
consists of $Q_i^2=fQ_i$ and the two monic length-eleven
relations, interchanged by $1\leftrightarrow2$.
All their coefficients belong to $R_0$, as the factorizations
in (P-F52f1)--(P-F52f3) verify. Replace a repeated adjacent
generator using the quadratic relation, and replace an
alternating length-eleven segment using its monic relation.
Every replacement strictly lowers word length and uses only
coefficients in $R_0$. Thus the $21$ words span the word
form over $R_0$; their independence over $F_0$ in the cited
presentation proves their independence over $R_0$ as well.

The reciprocal identities
\[
 \overline f=f,\qquad \overline L=L,\qquad
 \overline E_i=E_i,\qquad \overline N_0=N_0
                                                               \tag{P-F52aj}
\]
follow directly from (P-F52f1), where the bar on scalar
rational functions sends $q$ to $q^{-1}$.
Consequently all coefficients of the two monic relations
are reciprocal. The semilinear rule
\[
             \overline q=q^{-1},\qquad
             \overline {Q_i}=Q_i
                                                               \tag{P-F52ak}
\]
therefore defines an involutive algebra automorphism of
the actual word form and of $B_0$.
This verifies the generator-fixing internal bar directly from
the complete presentation. It is the coefficientwise word
convention used for the internal computations in GCT VIII,
Section~4, printed p.~27. We do not reconstruct the ambient
product-of-Hecke-algebras realization or its external
minimum-degree criterion here; the source distinguishes those
from the internal calculation in Sections~3.3--4 and p.~33.

The extension to the displayed splitting field is explicit:
\[
             D(q^{-1})=q^{-80}D(q),\qquad
             \overline x=q^{-40}x.
                                                               \tag{P-F52al}
\]
These equations preserve $x^2=D$ and square to the identity.
Furthermore $R(q^{-1})=q^{-40}R(q)$ in the notation of
(P-F52f1), so $g_1,g_2$ are bar-fixed. The other three
$g_i$ are reciprocal Laurent polynomials.
Writing $h_3(t)=t^3-E_1t^2+E_2t-E_3$, the source formulas
give
\[
 \begin{aligned}
 u_1^1&=(G-g_2)h_3(G)Q_1,
       &u_2^1&=(G-g_1)h_3(G)Q_1,\\
 v_0&=(2G-L)h_3(G)Q_1,
       &v_1&=-\frac{D}{q^{20}}h_3(G)Q_1.
 \end{aligned}                                                  \tag{P-F52am}
\]
In particular $u_1^1,u_2^1,v_0$ are bar-fixed,
$\overline {v_1}=q^{-40}v_1$, and $v_0$ belongs to
$\mathcal B_{\rm word}$.

The global element
\[
         \widehat u=FZ-a\,1\in Z(\mathcal B_{\rm word})
                                                               \tag{P-F52an}
\]
is integral in this word form and bar-fixed, because
$F,S,f,a$ are reciprocal Laurent polynomials with integer
coefficients. Its restriction to $W_0$ is $\mathsf u$.
Equation (P-F52ad) can now also be written $w=\widehat u v_0$,
so $w$ belongs to the same actual word form and is bar-fixed.
Equivalently, applying (P-F52al) to its explicit expression
in (P-F52ad) gives $\overline w=w$.
Thus
\[
          \mathcal L=R_0v_0\oplus R_0w
                     \subset\mathcal B_{\rm word}
                                                               \tag{P-F52ao}
\]
is an actual free rank-two submodule with bar-fixed basis,
stable under $FZ$ with the positive matrix (P-F52ae).
No saturation in the full word form or stability under the
individual $Q_i$ is asserted. Restricting
$R_0[\widehat u]$ to this submodule gives exactly the
faithful order $\mathcal O$ of (P-F52ac).
The quadratic relation is a relation of this restriction;
it is not asserted for $\widehat u$ on all of $B_0$.

\paragraph{A precise failure after the source's generator rescaling.}
The same global element need not belong to a generator form
after a nonunit rescaling. This can be decided here without
an ambient Hecke calculation. Let $C$ be either $R_0$, or
the source's coefficient ring $K_{\mathbb Q}$, the integral
closure of $\mathbb Q[q^{1/2},q^{-1/2}]$ in a splitting
field. This is GCT VIII, Section~2.2 p.~7, with the
$q^{1/2}$ convention specified in Section~3 p.~18.
For the latter choice, work after the corresponding scalar
extension of the actual algebra.
Consider the two precisely specified unital generator rings
\[
 \begin{split}
 \mathcal B_F&=C\langle FP_1,FP_2\rangle
              =C\langle FQ_1,FQ_2\rangle,\\
 \mathcal B_{\rm den}
      &=C\langle (f_p/f)P_1,(f_p/f)P_2\rangle.
 \end{split}                                                    \tag{P-F52ap}
\]
The first uses the centered positivity normalization
$F=q^{11}\widehat f_p$; a monomial centering does not
change its generated ring. The second is the literal
formal denominator rescaling of GCT VII, equation~(34),
p.~27. These rings are not being silently identified with
the convenient word form.

In every element of $\mathcal B_F$, the coefficient of
the basis word $Q_{12}$ is divisible by $F^2$ in $C$.
For a proof, start with a word of length $\ell$ in
$FQ_1,FQ_2$. It is $F^\ell$ times the corresponding
convenient word. Every monic reduction just described has
integral Laurent coefficients, so retains the scalar factor
$F^\ell$ and introduces no denominator. Starting lengths
$0$ and $1$ have no $Q_{12}$ coefficient, while every
starting length $\ell\ge2$ retains a factor divisible by
$F^2$. Taking finite $C$-linear combinations proves the
claim. The $21$ words remain a basis over $C$ because
they are independent over the splitting field.

In contrast, expansion of the actual global word gives
\[
 \widehat u
    =Ff(Q_1+Q_2)-F(Q_{12}+Q_{21})
                      +\bigl(F(S-f^2)-a\bigr)1.
                                                               \tag{P-F52aq}
\]
Its $Q_{12}$ coefficient is $-F$, which is not divisible
by $F^2$ in $C$. Indeed $F$ is a nonunit Laurent polynomial
by (P-F52s). If $F^{-1}$ were integral over
$\mathbb Q[q^{1/2},q^{-1/2}]$, a monic equation for it,
multiplied by the corresponding power of $F$, would give
$1\in F\mathbb Q[q^{1/2},q^{-1/2}]$, a contradiction.
This also excludes $F^{-1}\in R_0$.
Therefore
\[
                  \widehat u\notin\mathcal B_F.
                                                               \tag{P-F52ar}
\]
The exact source identity
$f_p/f=-q^{-12}(q-1)^2F$ from (P-F52r)--(P-F52s)
shows that $\mathcal B_{\rm den}\subset\mathcal B_F$,
so $\widehat u\notin\mathcal B_{\rm den}$ as well.
This is a membership assertion for the displayed generator
forms, not a claim that every source convention uses the
same normalization.

\paragraph{The positive basis is not a nonnegative change of the displayed source basis.}
There is also a concrete obstruction to a positive change
of basis. GCT VIII, Section~5.3 p.~31 and Figure~6 p.~38,
displays $u_1^1,u_2^1$ among its candidate canonical basis
elements. Figure~6 has exactly the same five coefficients
as GCT VII Figure~7 p.~36, including the lower coefficients
checked in (P-F52f1)--(P-F52f3). The source's identification
with its external minimum-degree construction is not
being promoted from the internal computation on p.~33.

The exact expansion of our second vector in this displayed
basis is
\[
        w=\left(\frac e2+\frac{Fx}{2q^{20}}\right)u_1^1
             +\left(\frac e2-\frac{Fx}{2q^{20}}\right)u_2^1.
                                                               \tag{P-F52as}
\]
For every real $q>0$, choose the branch $x=\sqrt{D(q)}>0$.
By P-F39, $D(q)>0$ and $F(q)>0$, and by (P-F52z),
$b(q)>0$. Equation (P-F52y) gives
\[
       \left(\frac{Fx}{2q^{20}}\right)^2
                          =\frac{e^2}{4}+b>\frac{e^2}{4}.
\]
Since $e(q)>0$, the second coefficient in (P-F52as) is
strictly negative and the first strictly positive.
The generic source-basis expansion is unique; thus this
is not a change of basis with coefficient functions
nonnegative on the positive real branch. This conclusion
concerns the generic coefficients, and does not assume
that the full source basis stays independent at every
exceptional specialization. It excludes the stated
nonnegative-cone comparison regardless of whether the
displayed candidate basis satisfies the source's remaining
canonicity conjectures. It is not a refutation of those
conjectures or of their relaxed positivity definitions.

The actual convenient word integrality and internal-bar
compatibility are therefore established, while the
specific rescaled-generator membership and nonnegative
change-of-basis comparisons just described fail.
An ambient-Hecke minimum-degree or categorical comparison
requires further structure beyond these finite properties.
\hfill$\square$

\paragraph{A structural sign limitation for a positive quadratic basis.}
The mixed spectral sign is unavoidable for any unit-containing
positive Laurent basis of a quadratic field. Precisely, let
$K'/\mathbb R(q)$ be a quadratic field, and suppose a
$\mathbb Z[q,q^{-1}]$-subalgebra is free on $1,u$, which
also form an $\mathbb R(q)$-basis of $K'$, with multiplication
\[
        u^2=\alpha(q)+\beta(q)u,
        \qquad
        \alpha,\beta\in\mathbb Z_{\ge0}[q,q^{-1}].
                                                               \tag{P-F52at}
\]
Then $\alpha\ne0$: otherwise $u(u-\beta)=0$ in the field,
forcing $u=0$ or $u=\beta$, contrary to the independence of
$1,u$. Consequently $\alpha(q)>0$ for every real $q>0$.
At each such specialization the two roots are real and have
opposite signs, since
\[
   u_\pm(q)=\frac{\beta(q)\pm
                     \sqrt{\beta(q)^2+4\alpha(q)}}2,
   \qquad u_-(q)<0<u_+(q).
                                                               \tag{P-F52au}
\]
Thus in the specialized regular representation, or any real
representation containing both characters, the element $u$
cannot act as a positive semidefinite operator for any
positive-definite inner product. A negative eigenvalue
alone rules this out. In the present order $\alpha=b$ and
$\beta=e$, so the formula agrees with the two eigenvalues
of $\mathsf u$. The assertion concerns simultaneous
nonnegative structure constants and positive operator
signs for a unit-containing field basis. A source canonical
basis is not required to consist of positive semidefinite
operators, and no such requirement is imposed here.
\hfill$\square$

\paragraph{A positive integral lattice for both actual rescaled generators.}
The quadratic order also controls a complete actual left module,
not only the central operator. Put
\[
 \widehat v=e\,1-\widehat u=d\,1-FZ,\qquad
 h=FS-d=\frac{FL-e}{2}.
                                                               \tag{P-F52av}
\]
The global element $\widehat v$ is central, integral in
$\mathcal B_{\rm word}$, and bar-fixed. Its restriction to
$W_0$ is $v=e-\mathsf u$, which satisfies
$v^2=ev+b$. Hence the same faithful quadratic order is
$\mathcal O=R_0[v]$, with its positive basis $1,v$.

Here $h$ is an integer Laurent polynomial with nonnegative,
symmetric, unimodal coefficients. Its complete central-to-edge
coefficient vector, in exponents $0,1,\ldots,24$, is
\[
 \begin{split}
 (164,164,160,156,152,149,143,137,125,114,104,94,83,\\
 \hspace{18mm}73,62,51,41,31,25,19,14,9,6,3,1).
 \end{split}                                                    \tag{P-F52aw}
\]
The coefficients in negative exponents are reflected, and
all remaining coefficients are zero. This is a complete
finite certificate: in (P-F52av) substitute (P-F52s) and
\[
       L=3(q^8+q^{-8})+4(q^4+q^{-4})
                        -2(q^2+q^{-2})+10;
\]
direct multiplication gives (P-F52aw). In particular
$h(1)=3996$.

Define four actual vectors of $B_0$ by
\[
  z_1=v_0,\qquad z_2=\widehat v v_0,\qquad
  z_3=Q_2v_0,\qquad z_4=\widehat v Q_2v_0.
                                                               \tag{P-F52ax}
\]
They are integral in $\mathcal B_{\rm word}$ and fixed by
the internal bar, by (P-F52ak)--(P-F52ao).
They are independent over $F_0$.
For a detailed proof, project after scalar extension to
the block $U_1$. The projections of $v_0,Q_2v_0$ are
$u_1^1,Q_2u_1^1$. First these two vectors are independent
over $K$. Indeed $Q_1u_1^1=fu_1^1$ because
$u_1^1=p_1(G)Q_1$. Also
\[
           Q_1Q_2u_1^1=Gu_1^1=g_1u_1^1.
                                                               \tag{P-F52ay}
\]
This follows from $GQ_1=Q_1Q_2Q_1$ and centrality.
If $Q_2u_1^1=\beta u_1^1$, the relation
$Q_2^2=fQ_2$ gives $\beta\in\{0,f\}$ and then
(P-F52ay) gives $g_1\in\{0,f^2\}$.
This is impossible: the nonzero coefficient of $x$ in
$g_1$ and the nonsquare $D$ show $g_1\notin F_0$.

On this block $\widehat v$ is the scalar
$r_1=d+Fc$, which likewise does not belong to $F_0$.
An $F_0$-linear relation among the four vectors in
(P-F52ax) therefore projects to
\[
       (a_1+a_2r_1)u_1^1+
                         (a_3+a_4r_1)Q_2u_1^1=0.
\]
Independence of the two displayed $K$-vectors, followed by
$r_1\notin F_0$, forces all four $a_i$ to vanish.
This proves the claimed independence without using the
quarantined generator matrices.

Let $\mathcal V=\bigoplus_{i=1}^4R_0z_i$.
The action of $\widehat v$ on its two pairs of vectors is
the same positive matrix
\[
             U=\begin{pmatrix}0&b\\1&e\end{pmatrix}.
                                                               \tag{P-F52az}
\]
Indeed its quadratic relation holds on $W_0$ and, by
centrality, also on $Q_2W_0$.
On the ordered basis (P-F52ax), actual left multiplication
by the two rescaled generators has the block matrices
\[
 \begin{split}
    \mathsf A_1&=
       \begin{pmatrix}FfI_2&hI_2+U\\0&0\end{pmatrix},\\
    \mathsf A_2&=
       \begin{pmatrix}0&0\\FI_2&FfI_2\end{pmatrix}.
 \end{split}                                                    \tag{P-F52ba}
\]
For verification, $Q_1v_0=fv_0$, $Q_2^2v_0=fQ_2v_0$,
and $Q_1Q_2v_0=Gv_0$ hold as above.
On this module
\[
                 FG=FS-FZ=h+\widehat v.
                                                               \tag{P-F52bb}
\]
Centrality extends the same formulas to $z_2,z_4$,
proving every column of (P-F52ba).
Every entry is in $\mathbb Z_{\ge0}[q,q^{-1}]$:
this has been certified for $b,e,h$, and is immediate
for $F$ and $f$. Thus $\mathcal V$ is a bar-fixed-basis
integral lattice stable under both $FQ_1$ and $FQ_2$.
Since $F$ is invertible over $F_0$, its scalar extension
is stable under $Q_1,Q_2$ and is an actual four-dimensional
left ideal of $B_0$. The displayed map from the free
rank-four module to these four source vectors is the
intertwiner, including its normalization. It is not an
identification of the earlier two-dimensional space
$\iota(K)$ with a left ideal.

In particular every ordered product of $FQ_1,FQ_2$, and
every nonnegative Laurent linear combination of such
products, has nonnegative Laurent matrix entries on this
lattice. Each of the two matrices also has an immediate
finite additive model: for every entry choose the finite
graded vector space with that character, and use its
tensor product and direct sum on four copies of finite
graded vector spaces. The two resulting additive
endofunctors have exactly (P-F52ba) on $K_0$; composition
models the matrices of all ordered words. The map
(P-F52ax), after scalar extension, intertwines every such
word with its actual source action. This does not supply
natural isomorphisms for the nonstandard defining
relations or a monoidal categorification of $B_0$.

There is a precise limitation on making both choices of
complementary generators positive in one basis. Suppose
a basis gave nonnegative real matrices for all four
$FQ_i,FP_i$ at a positive real parameter where the
representation remains the generic four-dimensional one.
Since
\[
             FQ_i+FP_i=FfI,\qquad P_i=f-Q_i,
                                                               \tag{P-F52bc}
\]
all off-diagonal entries of every matrix would vanish.
The $Q_i$ would then commute. They do not commute
generically on this module: if they did, applying
$Q_1Q_2=Q_2Q_1$ to $u_1^1$ would give
$g_1u_1^1=fQ_2u_1^1$, contradicting the independence
just proved. Consequently there is no common basis over
the function field whose four matrices are nonnegative
on an open interval of positive parameters with valid
specialization. Equivalently, at any such specialization
where the two actual $Q_i$ do not commute, a simultaneous
nonnegative basis for both complementary pairs is
impossible. This does not prevent the positive matrices
for the chosen pair in (P-F52ba), or assert that a source
canonical basis requires positivity for both pairs.
The new comparison concerns actual module actions;
the ambient canonical-basis, coherent categorical and
uniform-family questions remain additional requirements.
\hfill$\square$

\paragraph{The quadratic functor relations and the precise adjoint gap.}
The additive models just constructed can be chosen to respect
the quadratic generator relations by specified natural
isomorphisms. Let $\mathscr C$ be two copies of the category
of finite graded complex vector spaces. Choose graded spaces
$H,F_{\mathrm{gr}}$ with characters $Ff,F$, and let
$\mathscr C_h:\mathscr C\to\mathscr C$ be the matrix
functor with coefficient matrix $hI_2+U$. All its entries
have already been certified nonnegative. On
$\mathscr C\times\mathscr C$ take
\[
 \begin{split}
 \mathscr E_1(X,Y)&=(H\otimes X\oplus\mathscr C_h(Y),0),\\
 \mathscr E_2(X,Y)&=(0,F_{\mathrm{gr}}\otimes X\oplus H\otimes Y).
 \end{split}                                                    \tag{P-F52bd}
\]
Tensoring by a space is understood componentwise. These are
the actual matrices (P-F52ba) on $K_0$. Direct substitution
gives natural isomorphisms
\[
       \mathscr E_i^2\simeq H\otimes\mathscr E_i,
       \qquad
       \mathscr E_i^r\simeq
                     H^{\otimes(r-1)}\otimes\mathscr E_i
                     \quad(r\ge1).
                                                               \tag{P-F52be}
\]
For example the nonzero component of $\mathscr E_1^2(X,Y)$
is $H\otimes(H\otimes X\oplus\mathscr C_h(Y))$;
that of $\mathscr E_2^2(X,Y)$ is
$H\otimes(F_{\mathrm{gr}}\otimes X\oplus H\otimes Y)$.
The isomorphisms are precisely tensor associativity and
distributivity. Their repeated-run compatibilities follow
by their identical action on elementary tensors and the
specified direct-sum summands. This supplies coherent
reduction of a run of equal generators. It does not supply
natural isomorphisms or coherence for the mixed
length-eleven relation of the nonstandard algebra.

Ordinary adjunction in this explicit semisimple model does
not implement the source's self-adjoint involution.
For a matrix functor on finite graded vector spaces, its
left and right adjoints have coefficient matrix
\[
                         A(q^{-1})^{\mathsf T}.
                                                               \tag{P-F52bf}
\]
Indeed the adjoint of $V\otimes-$ is $V^*\otimes-$,
whose character is $\operatorname{char}(V)(q^{-1})$,
and the direct-sum input and output indices interchange.
The matrices in (P-F52ba) are reciprocal entrywise, but
neither is symmetric: each has a nonzero off-diagonal
block and a zero reverse block. Consequently neither
$\mathscr E_i$ is isomorphic to its own adjoint, even
up to an overall grading shift, because such an
isomorphism would give equal $K_0$ matrices and preserve
the zero pattern. This is an obstruction for the
specified direct-sum model with its ordinary adjunction.
It is not an obstruction for categories with different
Hom spaces or a different comparison with source objects.

\paragraph{Nonnegative Laurent Gram data exist for the actual matrices.}
In fact the matrix-level adjoint compatibility has an
explicit positive refinement. Set
\[
 \begin{split}
 T_0&=\begin{pmatrix}2&e\\e&e^2+2b\end{pmatrix},\\
 P_0&=T_0(hI_2+U),\qquad
 \mathsf H=\begin{pmatrix}FfT_0&P_0\\P_0&fP_0\end{pmatrix}.
 \end{split}                                                    \tag{P-F52bg}
\]
These are actual matrices with integer Laurent entries.
For completeness,
\[
 T_0U=\begin{pmatrix}
 e&e^2+2b\\ e^2+2b&e^3+3eb
 \end{pmatrix}.
                                                               \tag{P-F52bh}
\]
Thus every entry of $\mathsf H$ is nonnegative and
reciprocal, and $\mathsf H$ is symmetric. The identities
$T_0U=U^{\mathsf T}T_0$ and (P-F52ba) give directly
\[
             \mathsf H\mathsf A_i
                      =\mathsf A_i^{\mathsf T}\mathsf H
                       \qquad(i=1,2).
                                                               \tag{P-F52bi}
\]
For example the upper-right block in the first equality
is $FfT_0(hI_2+U)=FfP_0$ on both sides; the lower-right
block follows by the same self-adjointness of $U$ for
$T_0$. The second equality follows by direct block
multiplication. These are also the required numerical
graded-adjoint identities, since the coefficients are
reciprocal.

For every $q>0$ with $q\ne1$, the real matrix $\mathsf H(q)$
is positive definite. At $q=1$ it is positive semidefinite
of rank three. We verify these claims from the exact
source scalars, rather than assuming a categorical
realization. In the quadratic algebra generated by $v$,
the trace matrix of the pairing $(a,c)\mapsto
\operatorname{Tr}(ac)$ in the basis $1,v$ is $T_0$:
the two roots have sum $e$ and product $-b$.
They are distinct and real by (P-F52au).
The matrix $\mathsf H$ is consequently the trace of the
two-by-two cell form
\[
           \begin{pmatrix}
               Ff&h+v\\h+v&f(h+v)
           \end{pmatrix}.
                                                               \tag{P-F52bj}
\]
Under its two real embeddings, $h+v=Fg_i$ for
$i=1,2$. Each embedded form has positive leading entry
$Ff$ and determinant $F^2g_i(f^2-g_i)$.

Both $g_i$ lie in $(0,f^2)$ when $q>0$, $q\ne1$.
Their reality follows from $D(q)>0$ in P-F39; their
product is the strictly positive $N_0$ in (P-F52f1).
For their sum and the sum of their complements, put
$t=q^2+q^{-2}-2\ge0$. Exact expansion gives
\[
 \begin{split}
 L&=3t^4+24t^3+64t^2+62t+20,\\
 2f^2-L&=2t^5+17t^4+50t^3+64t^2+42t+12.
 \end{split}                                                    \tag{P-F52bk}
\]
Both are positive. The complement product is exactly
\[
 \begin{split}
 (f^2-g_1)(f^2-g_2)
 &=f^4-Lf^2+N_0\\
 &=\frac{(q^2-1)^4(q^4+q^2+1)(q^4-q^2+1)^2}{q^{20}}\,
       \left(\sum_{j=0}^4q^{2j}\right)
       \left(\sum_{j=0}^6q^{2j}\right).
 \end{split}                                                    \tag{P-F52bl}
\]
Every factor is positive for $q>0$, $q\ne1$.
The sum and product tests prove positivity of both
roots and both complements. Each embedded form in
(P-F52bj) is therefore positive definite; taking their
sum on the two distinct real embedding coordinates
proves the assertion for $\mathsf H$.
At $q=1$ one has $g_1,g_2=4,16$ and $f^2=16$.
One embedded form has rank two and the other rank one,
both positive semidefinite. The embedding-coordinate
change is still invertible, proving rank three.

This supplies actual nonnegative Laurent Gram data
compatible with both source-generator actions.
In particular the ordinary transpose-dual mismatch of
(P-F52bf) cannot be upgraded to a general adjunction
obstruction. A category using nonsimple objects or
nontrivial graded Hom spaces could have a different
Gram matrix. What remains is to realize the required
Hom compositions, natural adjunctions and mixed
nonstandard relation compatibly with the source
canonical data. The pairwise coefficient spaces and
identities (P-F52bg)--(P-F52bi) alone do not construct
those operations.\hfill$\square$
